% ------------------------------------------------------------
% Chapter 6: Conclusion
%
% Summarises and evaluates. Introduces no number that has not already appeared
% in Chapter 4 with a table behind it.
% ------------------------------------------------------------

\chapter{Conclusion}
\label{ch:conclusion}

This thesis asked whether large language model chatbots respond differently
when the user's age changes and nothing else does, and if so whether the
difference is the one a child-safety regime would want. Answering it required
an instrument that did not exist, since published safety evaluations vary the
request and hold the user constant. A benchmark of 200 scenarios crossed with
thirteen age conditions and three replicates was built, grounded in the duties
the Online Safety Act 2023 places on user-to-user services, and run against six
commercial models from five providers for 46,800 requests. Every reply was
scored against a frozen thirteen-label rubric by a classifier validated against
manual annotation with an admission threshold fixed in advance, and every effect
is paired within scenario and corrected within a declared family.

The answer to the first question is yes, and unambiguously. Refusal on
age-restricted requests is 45.5 percentage points higher for an explicit minor
age than an explicit adult age, and text written for an explicit minor age sits
1.77 grade levels lower. Both effects hold on all six models, both survive
correction, and neither is an artefact of the thresholds chosen to compute it.
Requests the benchmark expects to be answered at every age are almost never
refused, at a macro-average of 0.5\%, so the conditioning is not a general
increase in caution.

The answer to the second question is the contribution this thesis would defend.
The conditioning is real but uneven, in a pattern that an evaluation of either
dimension alone would miss. Safety conditioning is concentrated at the adult
boundary: half of the entire movement across fourteen years of age occurs in the
single step from seventeen to eighteen. Linguistic conditioning is spread across
the ladder, and the same step accounts for between 4.0 and 11.0\% of it. A model
can therefore be well calibrated to the statutory line and poorly calibrated to
the developmental range beneath it, which is what these six models are.
Conditioning is also about three times weaker when the age is implied by context
than when it is stated outright, and real users imply far more often than they
state.

Two secondary findings carry their own weight. Crossing the decision with the
delivered content, rather than reading a binary Refusal Rate, separates a
permissive model from one that agrees and then supplies nothing; on this panel
those are different models, and a Refusal Rate reports them identically. And
decomposing the reading-level change into the two terms of the formula shows it
is carried almost entirely by word choice rather than sentence structure, which
is the opposite of the intuitive account.

A third observation is exploratory. Every measure above reads a reply through a
closed vocabulary, so none can register a reply rewritten without changing
category, and comparing the replies themselves shows that this happens: an
explicit minor age separates the text of age-restricted replies by 0.189 where
the same change moves the control strata by roughly half as much. The measure
has no polarity, so it adds nothing to the direction of the safety finding. What
it adds is scope: age conditioning is wider than the categorical rates record.

Judged against its aims the approach has been largely successful, with two
qualifications. The instrument works: it isolates the age signal, it produces
effects large enough to be unambiguous, and its controls behave as controls
should. But the safety measurement stops at whether requested material was
delivered and does not adjudicate whether it was harmful, so the headline
directional rates are upper bounds rather than measurements of harm. And the
readability sample is not the response corpus, because a tenth of replies are
too short to measure and the loss is concentrated in particular models and
strata. Both limits are visible in the reported numbers rather than concealed by
them, and neither is resolved.

The most striking single result is the one that motivates the title. Across the
decade of childhood the ladder covers, from seven to seventeen, refusal on
age-restricted requests falls by twenty points. Across the single year from
seventeen to eighteen it falls by twenty-eight. These models have learned where
the law draws its line far more sharply than they have learned what a child is,
and the difference between those two things is the gap that a child-safety
regime built on age verification alone will not close.
